\begin{abstract}
非孤岛式微网的购电决策受到供需波动、预报更新与电价变化的共同影响。本文在统一的能量平衡和储能状态约束下，依次建立确定性调度、随机日前购电与日内滚动调整模型，以保障供电并降低购电费用。数据分析显示，供需存在日内错位、周周期与月度差异，为储能调度和情景构造提供依据。

针对问题一，考虑典型日供需与电价的时序错位，以全天购电费用最小为目标，建立\textbf{购电与储能联合线性规划模型}，协调储能的跨时段能量分配。求得最小日购电费用 $35118.60$ 元，较无储能基准\textbf{降低26.92\%}，并实现光伏全额消纳；进一步得到等价最优日初储电量区间 $[8550,8700]\,\mathrm{kWh}$，据此选取代表方案。

针对问题二，考虑负载与光伏的不确定性，建立\textbf{两阶段随机线性规划模型}，联合优化日前购电及事后储能调度与紧急购电。负载七日同刻相关系数为 $0.9722$，且周五、六日总负载偏低，据此采用周同期负载与净负荷日总量、前三日光伏日总量均值等特征进行\textbf{高斯核加权}。该特征组合采用平均能量分数进行效果验证，发现较等权基线降低 \textbf{58.706\%}。$2025$ 年 $2$ 至 $12$ 月结算费用为 \textbf{1699.85} 万元，紧急购电量占计划购电量的 \textbf{0.256\%}。为检验跨日近视假设，引入终端储能价值进行配对对照：在相同情景池、核权重及近视执行层下，远视模型较近视多支出 $38051$ 元，未达到预设 $0.5\%$ 节费采纳门槛，故\textbf{保留近视模型}。

针对问题三，引入日内光伏预报更新，采用\textbf{四阶段滚动随机规划与情景模型预测控制}，依据新预报重建剩余时段情景，冻结已执行决策并优化后续购电调整。基准配置采用 $0.8$ 分位购电下界与 $5$ 情景执行层，仅用 $0{:}00$ 预报的全年费用为 $1795.60$ 万元，在所比较的四种嵌套策略中最低；依次新增 $6{:}00$、$12{:}00$、$18{:}00$ 更新，费用分别增加 $112.53$ 万元、减少 $25.77$ 万元和 $4.72$ 万元，表明\textbf{增加预报更新并不必然降低费用}，其收益需与调整代价共同权衡。

针对问题四，将固定电价扩展为实时波动电价，基于历史信息进行多尺度电价预测，并构造\textbf{负载、光伏与电价的联合残差情景}，保留同期相关性。沿用上述基准配置，日前购电与日内滚动调整策略的全年结算费用分别为 $1862.76$ 万元和 $1939.87$ 万元。价格信息价值对比显示，基准配置下完美价格信息未降低回测费用；取消分位下界并将执行层情景数增至 $50$ 后，两种策略分别节费 $14.18$ 万元和 $16.15$ 万元，说明信息收益与计划保守程度及执行层情景覆盖有关。

\keywords{两阶段随机规划\quad 情景模型预测控制\quad 情景生成\quad 储能调度\quad 价格信息价值}
\end{abstract}